\section{Two semidirect constructions}\label{sec:app-twists}
The first construction, in source Section~4.1, is the join described in the common paper. It is flag-transitive and has parabolics $A_E\rtimes B_J$; its residual connectedness and firmness follow there. For thin inputs it is thin exactly when every old type orbit is a singleton. Consequently the literal claim that it is never thin is too strong. In particular, trivial type action gives the join of the original thin geometries. No faithfulness or normality of a Borel subgroup was used in the exact panel argument.

The second construction has different parabolics. Its identification with the common orbit construction is given next, including the thinness qualification that survives nonfaithful actions.

\subsection{Literal subgroups and the common orbit theorem}
Write $\alpha=(A,(A_x)_{x\in S})$ and
$\beta=(B,(B_i)_{i\in I})$ for flag-transitive finite-rank coset
geometries, with $A_E=\bigcap_{x\in E}A_x$ and
$B_J=\bigcap_{i\in J}B_i$. Empty intersections are the ambient
groups. Fix an indexed action of $B$ on $S$ together with an action
on $A$ satisfying $b(A_x)=A_{bx}$. This indexed action is necessary
data if distinct types have identical maximal parabolics; permutation
of the unlabelled subgroup set does not itself specify it.

In source Definition 4.5 there is one nonsingleton orbit $L$, a
representative $x_0\in L$, and singleton orbits elsewhere. Translate
its IPO convention into
\[
 B_Jx_0\cap B_Kx_0=B_{J\cup K}x_0
 \qquad(J,K\subseteq I).
\]
The same equation holds automatically on a singleton orbit. Choose
its unique element as representative. Let $\mathcal K$ be all these
orbits and set $T=I\sqcup\mathcal K$. For $J\subseteq T$ put
\[
 J_B=J\cap I,\quad J_A=J\cap\mathcal K,\qquad
 E(J)=S\setminus\bigcup_{M\in\mathcal K\setminus J_A}B_{J_B}x_M.
\]
The common orbit theorem gives the exact standard intersections
\begin{equation}\label{B-eq:intersection}
 G_J=A_{E(J)}\rtimes B_{J_B},\qquad
 G_T=A_S\rtimes B_I,
 \quad G=A\rtimes B.
\end{equation}
For a singleton orbit $M=\{y\}$, the union includes $y$ unless that
orbit type belongs to $J_A$. Consequently
$E(J)=(L\setminus B_{J_B}x_0)\cup\bigcup_{M\in J_A}M$.
At a maximal parabolic this is exactly Definition 4.6. Thus no
generated-support identity, or an interchange between quotient
images and intersections, is part of this identification.

The same theorem proves flag transitivity and chamber completion,
finiteness if and only if both inputs are finite, and preservation of
firmness and residual connectedness under their stated input
hypotheses. In particular it recovers Lemmas 4.7--4.8,
Propositions 4.9--4.11 and Theorem 4.12(a)--(d), with the corrected
$B_{J_B}$ notation in the second factor of Lemma 4.8. Its exact
panel indices are
\begin{align*}
 [G_{T\setminus\{i\}}:G_T]
   &=[B_{I\setminus\{i\}}:B_I] &&(i\in I),\\
 [G_{T\setminus\{M\}}:G_T]
   &=[A_{S\setminus B_Ix_M}:A_S] &&(M\in\mathcal K).
\end{align*}
These recover Theorem 4.12(e)'s firmness assertion. If the inputs are
thin and $B_I$ fixes every $x_M$, every displayed index is two.
In particular faithful regular input actions give the regular
conclusion of Corollary 4.13, with faithful full type-preserving
group $A\rtimes B$ when both chamber stabilizers are trivial.
The common theorem's chamber count is
$[G:G_T]=[A:A_S][B:B_I]$, including infinite cardinalities;
finiteness refers to the geometries, not to their possibly
unfaithful acting groups.

\subsection{Why thinness needs the chamber-stabilizer condition}
Here is the literal obstruction when thin coset models may act
unfaithfully. It is the same complete coordinate example already
used for the later source on wreath products; we give its N005
adapter explicitly. Take $A=C_2^3$ with the three coordinate
hyperplanes as its maximal parabolics. Its coset geometry is the
join of three two-point sets, a faithful regular rank-three geometry.
Let $B=S_3$ permute the coordinates and let
$\beta=(S_3,(C_3))$. This is a thin rank-one geometry with two
elements and nontrivial chamber stabilizer $C_3$.

There is one orbit $L$ of size three. Both $B$ and $C_3$ are
transitive on it, so the entire IPO condition holds, including the
empty and full subsets of the one-element type set. The source
twisting subgroups are, literally,
\[
 P_L=1\rtimes S_3,
 \qquad P_i=C_2^3\rtimes C_3.
\]
Their intersection is $1\rtimes C_3$. Every coset of $P_i$ meets
every coset of $P_L$: $P_iP_L=G$ because their factors contain the
whole base and the whole acting group. Hence the geometry is the
complete bipartite graph on parts of sizes two and eight. Its two
panel sizes are eight and two; it is firm, residually connected and
flag-transitive, but is not thin. Every assertion here follows
directly from the displayed finite groups, without a computation.

The source's proof of Theorem 4.12(e) replaces thinness by triviality
of both Borel subgroups. That implication requires faithfulness and
is false for the coset model above. Its thinness assertion and the
unqualified regular-input reading of Corollary 4.13 must therefore
retain the stated action hypothesis or the exact panel criterion.
This is a correction to literal scope, not a failure of the faithful
regular twisting theorem. No numbered-conjecture or historical
priority claim is attached to it.
